%%==================================================================FRONT PAGE AND TOC
%% For article only
%\mode<presentation:0>{\thispagestyle{empty}\maketitle}
%
%% For presentation only
%\mode<presentation| article:0| handout:0>{
%    \begin{frame}<article:0>[label=portada]
%    \titlepage
%    \end{frame}%Fin del frame
%}
%
%% For handout only
%\mode<handout>{
%  \begin{frame}[label=portada]
%    \maketitle
%  \end{frame}
%}
%
%%% TABLE OF CONTENTS
%\begin{frame}[label=toc]
%    \mode<article:0>{\frametitle{Contents}}
%    \mode<presentation>{\small}
%    \tableofcontents%[hidesubsections]
%\end{frame}
%
%%==================================================================S LOS DATOS
\section{Los datos}
%%==================================================================F 
\begin{frame}
\frametitle{La nube de puntos}
Metadatos
\begin{itemize}
 \item Sensor LiDAR: \alert{Optech ATLM}
 \item Datos disponibles: Primer y último impulso.
 \item Localización: Centro de Stuttgart (estación f.c.)
 \item Resolución: $0.91$ puntos$/m^2$ 
 \item Resolución media: $0.87$ puntos$/m^2$
 \item Numero de puntos: 
  \begin{enumerate}
    \item $259030$ de primeros retornos
    \item $259030$ de últimos retornos
   \end{enumerate}
\end{itemize}
\end{frame}
%%==================================================================F 
\pgfdeclareimage[width=0.5\textwidth]{stuttgart}{images/stuttgart_google}
\begin{frame}
\frametitle{Estación f.c. de Stuttgart}
 \begin{columns}
  \begin{column}{0.45\textwidth}
	Coordenadas UTM (MBR):
	\begin{itemize}
	 \item Norte = $5403764~m$
	 \item Sur = $5403188~m$
	 \item Oeste = $513632~m$
	 \item Este = $513118~m$
	\end{itemize}
   \end{column}
  \begin{column}{0.55\textwidth}
    \begin{center}
	\begin{tikzpicture}
    	  \pgftext[bottom,left,at={\pgfpointxy{0}{0}}]{\pgfuseimage{stuttgart}}
	\end{tikzpicture}
	\tiny{\textcopyright 2008 Google -- imágenes}
    \end{center}
   \end{column}
 \end{columns}
\end{frame}
%%==================================================================F 
\pgfdeclareimage[width=0.7\textwidth]{las}{images/las_format}
\begin{frame}
\frametitle{Estandard binario \emph{LAS} del ASPRS:}
\begin{itemize}
 \item Encabezado público
 \item Registros de longitud variable
 \item Puntos: X, Y, Z, Intensidad, clasificación, dirección de la pasada del vuelo,...
    \begin{center}
	\begin{tikzpicture}
    	  \pgftext[bottom,left,at={\pgfpointxy{0}{0}}]{\pgfuseimage{las}}
	\end{tikzpicture}
   \end{center}
 \item Para un conversor \alert{abierto} del estándar \emph{LAS}:
\beamergotobutton{\url{http://liblas.org/}}
\end{itemize}
\end{frame}
%%==================================================================F 
\defverbatim[colored]\head{
\begin{lstlisting}[style=shell]
  $ # Formato ASCII de los datos
  $ head csite4
\end{lstlisting}
}
\frame[containsverbatim]{%
\frametitle {Posible formato ASCII}
\head
\begin{center}
\begin{verbatim}
              X          Y       Z
          513628.82|5403190.04|291.40
          513628.85|5403191.47|291.24
          513628.95|5403192.93|294.31
          513628.99|5403194.45|294.17
          513628.97|5403196.05|291.26
          513629.01|5403197.58|291.24
          513629.05|5403199.10|291.25
          513629.09|5403200.53|291.26
                        ...
\end{verbatim}
\end{center}
}
%%==================================================================S
\section{Análisis LiDAR en GRASS}
%%==================================================================Sb
\begin{frame}
 \frametitle{Geographic Resources Analysis Support System\hfill 
   \includegraphics[width=0.7cm]{images/grasslogo_transp_big.png}} 
 \begin{itemize}
    \item Conocido como \alert<1>{GRASS}: GIS para manipulación y 
      análisis de datos geoespaciales.
    \item Como es un programa de \alert{códico abierto} con un 
      montón de librerias, es muy adecuado para la implementación 
      de nuestras propias herramientas.
    \item Escrito en C/C++, lo que permite implementar programas 
      computacionalmente pesados $\Rightarrow$ \alert{ideal} 
      para trabajar con LiDAR.
    \item GRASS $\geq$ 6 incluye indexación topológica $\Rightarrow$ 
      \alert{fatal} para trabajar con LiDAR.
    \item Hay versiones para GNU/Linux, MS-Windows, Mac-OSX, SUN, etc. 
      \alert{\textexclamdown\textexclamdown Todos lo pueden utilizar!!}
    \item Más información en: \beamergotobutton{\url{http://grass.osgeo.ogr}} 
  \end{itemize}
\end{frame}
%%==================================================================Sb
\subsection{Importando nube de puntos}
%%==================================================================F 
\defverbatim[colored]\extension{
\begin{lstlisting}[style=shell]
GRASS6: > # Primero conocemos la extension del archivo
GRASS6: > r.in.xyz -sg in=csite4
GRASS6: > # Fijo la region con la que trabajar
GRASS6: > g.region n= s= e= w=
GRASS6: > # Genero una visualizacion de las pasadas
GRASS6: > # considerando para cada celda el numero de puntos
GRASS6: > r.in.xyz in=csite4 out=pasadas method=n fs='|'
\end{lstlisting}
}
\begin{frame}
 \frametitle{Importar como raster: \LARGE\path{r.in.xyz}}
 \extension
\end{frame}
%%==================================================================F 
\defverbatim[colored]\xyz{
\begin{lstlisting}[style=shell]
# Generar un raster con la informacion LiDAR
r.in.xyz in=csite4 out=stuttgart_first fs='|' x=1 y=2 z=3
r.in.xyz in=csite4 out=stuttgart_last fs='|' x=4 y=5 z=6
\end{lstlisting}
}
\begin{frame}
 \frametitle{Solape de pasadas: \LARGE\path{r.in.xyz}}
 \xyz
 \pause
 \begin{center}
 \includegraphics[height=50mm]{images/rinxyz_first}~
 \includegraphics[height=50mm]{images/rinxyz_last}
 \end{center}
\end{frame}
%%==================================================================F 
\defverbatim[colored]\vascii{
\begin{lstlisting}[style=shell]
# Generar dos vectoriales con la informacion LiDAR
# de ultimo y primer impulso
v.in.ascii -tbzr in=csite4 out=stuttgart_first fs='|' x=1 y=2 z=3
v.in.ascii -tbzr in=csite4 out=stuttgart_last fs='|' x=4 y=5 z=6
\end{lstlisting}
}
\begin{frame}
 \frametitle{Importar como vectorial: \LARGE\path{v.in.ascii}}
 \vascii
 \pause
 \begin{center}
 \includegraphics[height=50mm]{images/vinascii_strips}
 \end{center}
\end{frame}
%%==================================================================F 
\defverbatim[colored]\vlidar{
\begin{lstlisting}[style=shell]
# Mostrar la informacion del LAS
v.in.lidar -p input=csite4.las
# Crear una nueva ``localizacion'' con la informacion del LAS
v.in.lidar -i input="Serpent Mound Model LAS Data.las" location=Stuttgart
# Cerrar y abrir GRASS en la nueva localizacion
# Importar el archivo LAS sin topologia ni tabla asociada
v.in.lidar -tb input="Serpent Mound Model LAS Data.las" output=
\end{lstlisting}
}
\begin{frame}
 \frametitle{Importar como .LAS: \LARGE\path{v.in.las}}
 \vlidar
 \pause
 \begin{center}
 \includegraphics[height=50mm]{images/vinascii_strips}
 \end{center}
\end{frame}
%%==================================================================F 
\begin{frame}
  \frametitle{Importar como LAS: \LARGE\path{v.in.las}}
 \begin{center}
 \includegraphics[height=50mm]{images/vinascii_strips}
 \end{center}
\end{frame}
%%==================================================================Sb
\subsection{Eliminando puntos groseros}
%%==================================================================F 
\defverbatim[colored]\voutlier{
\begin{lstlisting}[style=shell]
# Limpiamos el primer impulso
v.outlier in=stuttgart_first out=station_first \
> outlier=station_first_out soe=10 son=10 thres_o=25
# Limpiamos el ultimo impulso
v.outlier in=stuttgart_last out=station_last \
> outlier=station_last_out soe=10 son=10 thres_o=30
\end{lstlisting}
}
\begin{frame}[shrink=05]
 \frametitle{\path{v.outlier}}
\begin{beamerboxesrounded}[shadow=true]{\textbf{\path{v.lidar.edgedetection}}\texttt{ input=name output=name outlier=name [soe=value]
[son=value] [lambda\_i=value] [thres\_o=value]}}
\begin{itemize}
 \item \textbf{input}: name of the input vector map
 \item \textbf{output}: name of the output vector map
 \item \textbf{soe}: Interpolation spline step value in e-w direction
 \item \textbf{son}: Interpolation spline step value in n-s direction
 \item \textbf{lambda\_i}: Thychonov regularization weigth. (default: 0.1)
 \item \textbf{thres\_o}: Threshold for the outliers. (default: 50)
\end{itemize}
\end{beamerboxesrounded}
\voutlier
\end{frame}
%%==================================================================F 
\begin{frame}
 \frametitle{Visualización de \LARGE\path{v.outlier}}
 \begin{center}
 \includegraphics[height=50mm]{images/voutlier_first}~
 \includegraphics[height=50mm]{images/routlier_first}
 \end{center}
\end{frame}
%%==================================================================F 
\begin{frame}
 \frametitle{Visualización de \LARGE\path{v.outlier}}
 \begin{center}
 \includegraphics[height=50mm]{images/voutlier_last}~
 \includegraphics[height=50mm]{images/routlier_last}
 \end{center}
\end{frame}
%%==================================================================Sb
\subsection{Detección de borde de objetos}
%%==================================================================F 
\defverbatim[colored]\vedge{
\begin{lstlisting}[style=shell]
# Deteccion de bordes
v.lidar.edgedetection in=station_last see=4 sen=4 lambda_g=0.1 \
> out=station_edge
\end{lstlisting}
}
\begin{frame}[shrink=02]
 \frametitle{\path{v.lidar.edgedetection}}
\begin{beamerboxesrounded}[shadow=true]{\textbf{\path{v.lidar.edgedetection}}\texttt{ input=name output=name first=name see=value sen=value [lambda\_g=value] [tgh=value] [tgl=value] [theta\_g=value] [lambda\_r=value]}}
\begin{itemize}
 \item \textbf{input}: name of the input vector map
 \item \textbf{output}: name of the output vector map
 \item \textbf{see}: Interpolation spline step value in e-w direction
 \item \textbf{sen}: Interpolation spline step value in n-s direction
 \item \textbf{lambda\_g}: Regularization weight in gradient evaluation (default: 0.01)
 \item \textbf{tgh}: High gradient threshold for edge classification (default: 6)
 \item \textbf{tgl}: Low gradient threshold for edge classification (default: 3)
 \item \textbf{theta\_g}: Angle range for same direction detection (default: 0.26)
 \item \textbf{lambda\_r}: Regularization weight in residual evaluation (default: 2)
\end{itemize}
\end{beamerboxesrounded}
\vedge
\end{frame}
%%==================================================================F 
\pgfdeclareimage[width=0.55\textwidth]{edge}{images/edge}
\begin{frame}
 \frametitle{Visualización de \LARGE\path{v.lidar.edgedetection}}
\begin{columns}
  \begin{column}{0.7\textwidth}
    \begin{center}
	\begin{tikzpicture}
    	  \pgftext[bottom,left,at={\pgfpointxy{0}{0}}]{\pgfuseimage{edge}}
	\end{tikzpicture}
   \end{center}
  \end{column}
  \begin{column}{0.25\textwidth}
	\begin{enumerate}
	 \item \textcolor{red}{Borde}
	 \item \textcolor{yellow!95!black}{Terreno}
	\end{enumerate}
  \end{column}
\end{columns}
\end{frame}
%%==================================================================Sb
\subsection{Relleno de los límites de los objetos}
%%==================================================================F 
\defverbatim[colored]\vgrow{
\begin{lstlisting}[style=shell]
# IMPORTANTE: Disminuir primero la resolucion
g.region -p res=2
v.lidar.growing in=station_edge out=station_grow \
> first=station_first
\end{lstlisting}
}
\begin{frame}[shrink=03]
 \frametitle{\path{v.lidar.growing}}
\begin{beamerboxesrounded}[shadow=true]{\textbf{\path{v.lidar.growing}}\texttt{ input=name output=name first=name [tj=value] [td=value]}}
\begin{itemize}
 \item \textbf{input}: name of the input vector map
 \item \textbf{output}: name of the output vector map
 \item \textbf{first}: name of the input vector map of the first pulse points
 \item \textbf{tj}: Threshold in planimetric units for considering two measurements as corresponding first and double pulse. (default
0.2)
 \item td: Threshold in height units for classifying a point as double pulse. (default: 0.6)
\end{itemize}
\end{beamerboxesrounded}
\vgrow
\end{frame}
%%==================================================================F 
\pgfdeclareimage[width=0.55\textwidth]{grow}{images/grow}
\begin{frame}
 \frametitle{Visualización de \LARGE\path{v.lidar.growing}}
\begin{columns}
  \begin{column}{0.65\textwidth}
    \begin{center}
	\begin{tikzpicture}
    	  \pgftext[bottom,left,at={\pgfpointxy{0}{0}}]{\pgfuseimage{grow}}
	\end{tikzpicture}
   \end{center}
  \end{column}
  \begin{column}{0.35\textwidth}
	\begin{enumerate}
	 \item \textcolor{yellow!95!black}{Terreno}
	 \item \textcolor{green}{Terreno doble eco}
	 \item \textcolor{blue!90!black}{Objeto doble eco}
	 \item \textcolor{red}{Objeto}
	\end{enumerate}
  \end{column}
\end{columns}
\end{frame}
%%==================================================================Sb
\subsection{Corrección final}
%%==================================================================F 
\defverbatim[colored]\vcorr{
\begin{lstlisting}[style=shell]
v.lidar.correction in=station_grow sce=60 scn=60 lambda_c=2 \
> tcl=0.1 out=station_corr terrain=station_terr
\end{lstlisting}
}
\begin{frame}[shrink=04]
 \frametitle{\path{v.lidar.correction}}
\begin{beamerboxesrounded}[shadow=true]{\textbf{\path{v.lidar.correction}}\texttt{ input=name  output=name  terrain=name  sce=value  scn=value  [lambda\_c=value]  [tch=value]  [tcl=value]}}
\begin{itemize}
 \item \textbf{input}: name of the input vector map
 \item \textbf{output}: name of the output vector map
 \item \textbf{terrain}: name of the terrain output vector map
 \item \textbf{sce}: Interpolation spline step value in e-w direction
 \item \textbf{scn}: Interpolation spline step value in n-s direction
 \item \textbf{lambda\_c}: Regularization weight in reclassification evaluation. (default 1)
 \item \textbf{tch}: High difference threshold (terrain/object). (default 2)
 \item \textbf{tcl}: Low differnce threshold (object/terrain). (default 1)
\end{itemize}
\end{beamerboxesrounded}
\vcorr
\end{frame}
%%==================================================================F 
\pgfdeclareimage[width=0.55\textwidth]{corr}{images/correction}
\begin{frame}
 \frametitle{Visualización de \LARGE\path{v.lidar.correction}}
\begin{columns}
  \begin{column}{0.65\textwidth}
    \begin{center}
	\begin{tikzpicture}
    	  \pgftext[bottom,left,at={\pgfpointxy{0}{0}}]{\pgfuseimage{corr}}
	\end{tikzpicture}
   \end{center}
  \end{column}
  \begin{column}{0.35\textwidth}
	\begin{enumerate}
	 \item \textcolor{yellow!95!black}{Terreno}
	 \item \textcolor{green}{Terreno doble eco}
	 \item \textcolor{blue!90!black}{Objeto doble eco}
	 \item \textcolor{red}{Objeto}
	\end{enumerate}
  \end{column}
\end{columns}
\end{frame}
%%==================================================================Portada
%\repeatframe{portada}
